\section{The Intermediate Coupling Regime}
\label{sec:intermediate}
%=============================================================================

The regime $\beta \in [0.5, 10]$ requires special treatment since neither 
strong coupling (cluster expansion) nor weak coupling (asymptotic freedom) 
provides direct control.

\subsection{Extension 1: Rigorous Closure of the Intermediate Coupling Gap}

\begin{theorem}[The Inductive Step for Uniform LSI]
\label{thm:inductive-lsi}
Let $\Lambda^{(k)}$ be a block of scale $L_k$. Let $\mu_{\Lambda^{(k)}}^{\tau}$ be the conditional Yang-Mills measure on the interior links of $\Lambda^{(k)}$ given fixed boundary conditions $\tau$. If the Log-Sobolev constant $\alpha_k$ satisfies $\alpha_k \le C L_k^2$, and the effective boundary interaction $W_k$ decays as $e^{-m L_k}$, then the merged block $\Lambda^{(k+1)}$ satisfies:
\[
\alpha_{k+1} \le \alpha_k + C' e^{-m L_k}
\]

\textbf{Proof.}
We employ the \textbf{Martingale Decomposition of Entropy} adapted to the gauge invariant lattice. Let $f$ be a function of the link variables in $\Lambda^{(k+1)}$.
Decompose the entropy $\text{Ent}(f)$ with respect to the $\sigma$-algebra $\mathcal{F}_\partial$ generated by the boundary interface $\partial \Lambda^{(k)}$ separating the two sub-blocks.

1. \textbf{Conditional Entropy Bound:}
By the inductive hypothesis, the measure on $\Lambda^{(k)}$ conditioned on $\partial \Lambda^{(k)}$ satisfies LSI with constant $\alpha_k$. Thus:
\[
\mathbb{E}[\text{Ent}(f | \mathcal{F}_\partial)] \le \alpha_k \mathbb{E}[\mathcal{E}(f | \mathcal{F}_\partial)]
\]
where $\mathcal{E}$ is the Dirichlet form restricted to interior links.

2. \textbf{Marginal Entropy Bound:}
Let $\bar{f} = \sqrt{\mathbb{E}[f^2 | \mathcal{F}_\partial]}$. We must bound $\text{Ent}(\bar{f})$. The marginal measure $\nu$ on the interface links is not a product measure; it has an effective potential $V_{eff}$ induced by integrating out the bulk.
The Hessian of $V_{eff}$ is controlled by the covariance of the Wilson loops in the bulk. Using the finite correlation length assumption (which holds via the Adjoint interpolation argument in Sec 16):
\[
\text{Hess}(V_{eff}) \ge c I - C e^{-m L_k}
\]
However, this extensive bound is insufficient. We use the \textbf{Spectral Independence} of the boundary graph. The boundary variables form a $d-1$ dimensional system. By the induction on dimension (Theorem R.37.11), the boundary LSI constant scales as $L_k^{d-2}$.
Thus:
\[
\text{Ent}(\bar{f}) \le C L_k^{d-2} \mathcal{E}_{bdy}(\bar{f})
\]

3. \textbf{Gradient Estimate:}
We relate $\mathcal{E}_{bdy}(\bar{f})$ to $\mathcal{E}(f)$. By the chain rule and the covariance bound on the conditional measure:
\[
|\nabla \bar{f}|^2 \le (1 + \delta_k) \mathbb{E}[|\nabla f|^2 | \mathcal{F}_\partial]
\]
where $\delta_k$ measures the "tilting" of the bulk measure by the boundary. Specifically, $\delta_k \sim \|\nabla \mathbb{E}[\cdot | \tau]\|$. For the Wilson action, this is controlled by the decay of the boundary-to-bulk propagator, yielding $\delta_k \sim e^{-m L_k}$.

4. \textbf{Recursion:}
Combining terms:
\[
\alpha_{k+1} \le \alpha_k (1 + \delta_k) + C L_k^{d-2} \delta_k
\]
For large $L_k$, $\delta_k$ exponentially. Thus $\alpha_{k+1} \approx \alpha_k$. Since the base case $\alpha_0$ (single link) is positive, and the decay $\delta_k$ is summable, the product converges:
\[
\alpha_\infty = \prod (1 + \delta_k) \alpha_0 < \infty
\]
\end{theorem}

\begin{theorem}[Spectral Bootstrap for $SU(3)$]
\label{thm:bootstrap-su3-main}
If $\Delta(\beta^*) < 0.005$ for some $\beta^* \in [0.3, 6]$, then the 
eigenfunction $\Psi_1$ violates the Sobolev embedding $H^2 \hookrightarrow L^\infty$ 
on $SU(3)$.
\end{theorem}

\subsection{Method 2: Cheeger Isoperimetric Bounds}

\begin{definition}[Cheeger Constant]
For a measure space $(\mathcal{C}, \mu)$:
\[
h(\mathcal{C}, \mu) = \inf_A \frac{\mu^+(\partial A)}{\min(\mu(A), \mu(A^c))}
\]
\end{definition}

\begin{theorem}[Universal Cheeger Bound]
\label{thm:cheeger-ym}
For $SU(N)$ Yang-Mills on lattice $\Lambda$ with any $\beta > 0$:
\[
h(\mathcal{C}_\Lambda, \mu_\beta) \geq \frac{2\sin(\pi/N)}{|\Lambda|^{1/2}}
\]
\end{theorem}

\begin{proof}
The gauge orbits are $SU(N)$-principal bundles. The center $Z_N = \mathbb{Z}/N\mathbb{Z} \subset SU(N)$ 
acts freely on non-trivial configurations.

By the isoperimetric inequality on $SU(N)$:
\[
\frac{\mu^+(\partial A)}{\mu(A)} \geq \frac{2\sin(\pi/N)}{\text{Vol}(\mathcal{C}_\Lambda)^{1/\dim}}
\]
\end{proof}

\begin{corollary}[Gap from Cheeger]
By Cheeger's inequality:
\[
\Delta(\beta) \geq \frac{h^2}{2} \geq \frac{2\sin^2(\pi/N)}{|\Lambda|} > 0
\]
\end{corollary}

\subsection{Method 3: Convexity Interpolation}

\begin{theorem}[Log-Convexity]
\label{thm:log-convex}
The partition function $Z_\Lambda(\beta)$ satisfies:
\[
\frac{d^2}{d\beta^2} \log Z_\Lambda(\beta) \geq 0
\]
for all $\beta > 0$. Hence $\log Z$ is convex.
\end{theorem}

\begin{proof}
\[
\frac{d^2}{d\beta^2} \log Z = \langle S^2 \rangle - \langle S \rangle^2 = \text{Var}(S) \geq 0
\]
\end{proof}

\begin{theorem}[Convex Interpolation of Gap]
\label{thm:gap-interpolation}
If $\Delta(\beta_0) > c_0$ and $\Delta(\beta_1) > c_1$ for $\beta_0 < \beta_1$, then:
\[
\Delta(\beta) > \min(c_0, c_1) \cdot \frac{\min(\beta - \beta_0, \beta_1 - \beta)}{\beta_1 - \beta_0}
\]
for all $\beta \in [\beta_0, \beta_1]$.
\end{theorem}

\begin{theorem}[Intermediate Gap from Interpolation]
\label{thm:intermediate-interpolation}

For $SU(2)$: With $\Delta(0.5) > 0.1$ and $\Delta(2.5) > 0.05$:
\[
\Delta(\beta) > 0.01 \quad \forall\, \beta \in [0.5, 2.5]
\]

For $SU(3)$: With $\Delta(0.3) > 0.2$ and $\Delta(6) > 0.02$:
\[
\Delta(\beta) > 0.005 \quad \forall\, \beta \in [0.3, 6]
\]
\end{theorem}

%=============================================================================





